% Pursuing Stacks (cotes 134-2 to 134-8): archive cote 134-6, site page 24 (PDF page 25), typescript page 343.
% typescript. Draft compared: ps5.tex:2552-2613.
% AI-transcribed by Claude Opus 5.5 (claude-opus-5-5, adaptive thinking, effort high), Message Batches API, batch ps3 (msgbatch_01XdFmuVVBAcgc7EFxdtCqsj), 2026-11-12; not yet checked by eye.
\scanpage{134-6/24}
7.7. \mspage{343}

95. After this cascade of ``afterthoughts'' on abelianization of homo-%
topy types, it is time now to resume some more technical work\hand{,} and get
through with this unending part IV, in accordance with the short range
working program I had come to four days ago (section 92, p. 320). I'll
take up the three topics stated there - namely contractors, induced
structures, and ``miscellaneous'' - in \xout{\illegible{1}}\inserted[typed]{that} order, as reviewed previous-%
ly. Thus we'll start with contractors. I have in mind now mainly the
definition of contractors, and a few basic facts following easily from
what is \xout{now}\inserted[typed]{already} known to us.

\deleted{\xout{Proposition}\inserted[typed]{Definition} 1 . Let A be a category. We say A is a \uline{precontractor} if
there exists on A a contractibility structure (section 51, D) such
that all elements of A are contractible.}

\xout{P} The first thought that comes to my mind is to define a contractor
as a category A such that the set \xout{of} Ob(A) of all objects of A is a con-%
tractibility structure on A\hand{,} \xout{(section 51, D)} i.e. that there exists
\xout{such a struct} a contractibility structure on A for which every object in
A is contractible. The trouble with this definition is that it makes
the implicit assumption that A is stable under finite products - as the
notion of a contractibility structure was defined only in a category
satisfying this extra assumption (cf. section 51, D)). Now this assump-%
tion is \uline{not} satisfied by the three standard test categories, including
\hand{$\Delta$}, which surely we do want to consider as contractors ! The next
thought then, suggested by this reflexion, is to embed A into $A^{\wedge}$ to
supply the products which may be lacking in A, and demand there exist
a contractibility structure on $A^{\wedge}$, such that the objects in A \xout{are} be
contractible\inserted{and moreover generate}; or, what amounts to the same, that \inserted{for} the homotopy
interval structure on $A^{\wedge}$ defined by intervals coming from A as a gene-%
rating family, \xout{be a contractibility structure admitting} the objects of
A are contractible (which implies that this structure ``is'' indeed a
contractibility structure). This condition (in the more general case\unsure{\hand{,}}
when A appears as a full subcategory of any larger category \unsure{\hand{M}}\inserted{(stable under finite products,}) has
been restated in wholly explicit terms as the ``\uline{basic assumption}''
(Bas 4) on a set of objects, in order that it generate a contractibi-%
lity structure (section 51, p.118). It is immediate that in the case
when A itself is stable under finite products in the ambiant category,
that this condition is intrinsic to A and just \xout{\illegible{1}} amounts to the
first definition we had in mind.

\margin{24}
